
@@M0@@
\item
\emph{@@M1@@ est num\'eriquement \'equivalent \`a @@M2@@.}
\end{enumerate}
\emph{Si @@M3@@ est projectif et muni d'un module inversible ample @@M4@@, les conditions pr\'ec\'edentes sont aussi \'equivalentes aux suivantes}~:
\begin{enumerate}[label=(\roman*),start=4]
\item \emph{Pour tout entier @@M5@@ (positif ou n\'egatif), @@M6@@ est ample.}
\item \emph{(Si @@M7@@ est irr\'eductible). Pour tout couple d'entiers @@M8@@, on a}
@@M9@@
\item \emph{(Si @@M10@@ est irr\'eductible de dimension @@M11@@). On a}
@@M12@@
\end{enumerate}

D\'emontrons que (i) implique (ii). Appliquant le lemme de d\'evissage (EGA III 3.1) et 4.5 (i), on se ram\`ene au cas o\`u @@M13@@ est int\`egre et @@M14@@. Alors on doit prouver que le polyn\^ome @@M15@@ est constant. Pour ceci, on peut remplacer @@M16@@ par un multiple, donc supposer @@M17@@ alg\'ebriquement \'equivalent \`a @@M18@@. Dans ce cas, si @@M19@@ est d\'eform\'e en @@M20@@ par la famille alg\'ebrique d\'efinie par un faisceau inversible @@M21@@ sur un @@M22@@ o\`u @@M23@@ est connexe, de type fini, le polyn\^ome @@M24@@ ne d\'epend pas de @@M25@@ (EGA III 7.9.5)\,; d'o\`u l'implication.

Soit @@M26@@ une courbe int\`egre ferm\'ee dans @@M27@@. Faisant @@M28@@ dans (ii), on trouve @@M29@@. Donc la formule (``th\'eor\`eme de Riemann'') @@M30@@ donne @@M31@@\,; d'o\`u (ii) implique (iii).

Soit @@M32@@ une pr\'esentation de Chow~: @@M33@@ est projectif et @@M34@@ est surjectif. Si l'on admet que (iii) implique (i) pour @@M35@@, alors 
